\section{Method}
We used a mixed-methods design that combined participant observation, video-based conversation analysis, and reflexive thematic analysis to examine ability-diverse collaboration in accessible video creation workshops. 

\subsection{Data Collection}

Data collection occurred across seven workshops conducted between April and July 2025, employing multiple methods to capture the complexity of ability-diverse interactions:

\textbf{Participant Observation}: The first author served dual roles as workshop facilitator and participant-observer, documenting interactions through detailed field notes taken during and immediately after each 3–4-hour session. This insider perspective enabled nuanced understanding of both planned activities and emergent collaborative moments. Field notes focused on: (1) instances of cross-ability collaboration, (2) technology-mediated interactions, (3) accessibility barriers and workarounds.

\textbf{Video Documentation}\textbf{:} All workshops were video recorded with participant consent with prior written consent obtained before the first workshop. Recordings captured multiple perspectives including screen recordings of digital activities, wide-angle views of group interactions, and close-ups of specific technical demonstrations. This visual data proved essential for analysing non-verbal communication, spatial arrangements, and the role of sighted guides in facilitating access.

\textbf{Post-Workshop Reflections}\textbf{:} Each workshop concluded with a 20–30-minute guided reflection session where participants discussed their learning experiences, challenges encountered, and collaborative strategies developed. These sessions were audio recorded and transcribed, providing immediate participant perspectives on the workshop activities.

\textbf{Participant Surveys}\textbf{:} Following completion of all workshops, participants completed an online survey (n=10) contains open-ended questions assessing challenges encountered and experiences of collaboration. Survey design followed accessibility guidelines, available in both screen-reader compatible and large-print formats, as well as offered in the form of telephone/Zoom survey.

\textbf{Semi-Structured Interviews}\textbf{:} Two BVI tutors participated in 60–90-minute semi-structured interviews exploring their experiences of the teaching process and cross-ability collaboration. Interview topics included: pedagogical adaptations for mixed-ability groups, observed collaboration patterns, technology's mediating role, and recommendations for future workshops. Interviews were conducted in 1 month after workshop completion, allowing for reflective distance.



\subsection{Data Analysis}

We employed a multi-layered analytical approach suited to our diverse data types:

\textbf{Video-Conversation Analysis} We employed Multimodal Ethnomethodological Conversation Analysis (EMCA) following Due et al. 's work \cite{due_distributed_2021}. The first author manually transcribed selected video clips, then used ATLAS.ti \cite{atlasti_scientific_software_development_gmbh_atlasti_2023} to code multimodal interactions capturing both verbal and embodied elements. This approach helped identify "collaboration episodes" - bounded sequences where participants worked together to achieve specific video creation goals. Each episode was coded for: (1) initiator (BVI participant or sighted instructor), (2) collaboration type, and (3) interaction modality (verbal, visual, tactile).

\textbf{Thematic Analysis} Post-workshop reflections, surveys, and semi-structured interviews were analysed thematically following Braun and Clarke's six-phase approach \cite{braun_reflecting_2019}. Initial codes were generated inductively from participants' descriptions of barriers, facilitators, and collaborative experiences. These codes were then refined through iterative review, resulting in broader themes about accessibility challenges, learning strategies, and support needs.

Following initial separate analyses, we conducted integration sessions where findings from different data sources were synthesized. This process involved mapping collaboration patterns identified through video analysis against themes from interviews and reflections, using interdependence theory as an organizing framework. Contradictions between data sources were explicitly noted and explored as potential insights into the complexity of ability-diverse collaboration.

\textbf{Coding Process}

All initial coding was conducted independently by the first author, then discussed with the research team during regular meeting sessions to review coded excerpts, discuss disagreements, and refine code definitions. Discrepancies were resolved through consensus, ensuring analytical rigor and consistency across the dataset.

\textbf{Integration of Findings}

The video-based EMCA findings and thematic analysis were integrated to produce a comprehensive understanding of BVI video creation education. Collaboration episodes identified through EMCA were contextualized using participants' reflective accounts, while survey responses validated observed interaction patterns. This triangulation revealed how micro-level collaborative practices connected to broader structural barriers in video creation education.

\subsection{Ethical Considerations}
This study received institutional ethics approval [PROTOCOL NUMBER PLACEHOLDER]. Given the visual nature of video workshops and varying vision abilities among participants, we developed a multi-modal consent process including verbal explanation, printed documents, and screen-reader accessible digital forms. Participants could choose their level of data sharing, with options to exclude video footage while maintaining audio participation.
Data management followed institutional guidelines for sensitive data, with video files stored on encrypted servers and access limited to the research team. All identifying information was removed from transcripts, and participants chose pseudonyms for publication.

